% =============================================================================
% background.tex -- Background and key concepts.
%
% REVISION CONTRACT (advisor feedback, 2026-07):
% * Section 2.4 (sec:concepts) defines EVERY technical term and statistical tool the
%   paper uses, each as "what it is / why it is needed / how it serves the research
%   question". Nothing later in the paper may use a term that is not defined here.
% * Section 2.5 (sec:chain) replaces the previous one-sentence assertion that
%   contamination indicates leakage with four explicit assumptions A1-A4, and marks
%   which the paper tests and which it does not. This is the conceptual spine.
% =============================================================================

\section{Background and Key Concepts}
\label{sec:background}

\subsection{Benchmarks as proxies for capability}
LLM benchmarks act as \emph{proxies} for abilities that cannot be measured directly, such
as reasoning, comprehension, factual recall, and programming skill. By scoring a model on
a fixed set of standardized tasks, the community infers how useful, and increasingly how
safe, that model is likely to be in deployment~\citep{hendrycks2021mmlu}. Familiar
examples target distinct competencies: MMLU covers multitask knowledge across 57
subjects~\citep{hendrycks2021mmlu}, GSM8K covers multi-step arithmetic word
problems~\citep{cobbe2021gsm8k}, and HumanEval covers functional code
generation~\citep{chen2021humaneval}. Scores on these suites drive model selection,
leaderboard position, and public claims of progress.

\subsection{The validity assumption, and what breaks when it fails}
The inference from benchmark score to capability rests on one strict assumption:
\emph{the test data was not seen during training}. Only then does a high score support
the intended conclusion, that the model \emph{generalizes}, meaning it applies learned
regularities to inputs it has not encountered, rather than \emph{memorizes}, meaning it
reproduces specific training instances.

When the assumption fails, the benchmark stops measuring what it claims to measure. A
memorized test item raises the score with no corresponding gain in ability, so the metric
becomes an unreliable estimate of the underlying construct. This distinction is not only
conceptual. Memorization is directly measurable as verbatim regeneration of training
text, and it grows predictably with model size, with how often a sequence was duplicated
in training, and with how much context the model is
given~\citep{carlini2023quantifying}. A closely related privacy-facing measurement plants
a secret in the training data and records how strongly the model prefers it over random
alternatives, which rises with the number of times the secret was
seen~\citep{carlini2019secret}. Memorization is also highly uneven across examples rather
than a uniform background rate, so which particular items a model memorizes is itself a
measurable per-example property~\citep{zhang2023counterfactual}.

\subsection{Static test sets meet weakly filtered corpora}
The structural tension is that benchmarks and training corpora have opposite properties.
Benchmarks are static, small, widely circulated, and publicly indexed. Once an MMLU or
GSM8K item is published it is copied into papers, blog posts, code repositories, and
forum discussions. Training corpora are the reverse: very large web scrapes with light
filtering. Common Crawl~\citep{commoncrawl} and The Pile~\citep{gao2020pile} are
assembled at the scale of hundreds of gigabytes to petabytes, where reliably removing any
particular short string is impractical.

The consequence is that benchmark items enter training corpora through ordinary
redistribution, with no adversary required. A public benchmark is therefore a persistent
and low-effort exposure surface, because the very properties that make a benchmark useful
(stable, shared, citable) are what make its eventual presence in a future crawl likely.

\subsection{Key concepts and methodological tools}
\label{sec:concepts}
This subsection defines every technical idea used later in the paper. Each entry states
what the concept is, why the study needs it, and how it bears on the research questions
of Section~\ref{sec:intro}.

\subsubsection{Token probabilities, loss, and perplexity}
A language model reads text as a sequence of \emph{tokens}, which are words or word
fragments. At each position it produces a probability distribution over which token comes
next. Given a real text, we can therefore ask what probability the model assigned to each
token that actually occurred.

The \emph{loss} of a text is the average negative logarithm of those probabilities. Low
loss means the model found the text unsurprising. \emph{Perplexity} is the exponential of
the loss and carries the same information in a different scale, so we use loss
throughout. Every detector in this paper is computed from these same per-token
probabilities, which is the fact that ultimately drives our result.

\emph{Why this is needed.} Loss is the oldest and simplest membership signal, on the
reasoning that a model tends to find its own training data less surprising than new
text~\citep{yeom2018privacy}. It is also the natural baseline against which any more
elaborate detector must prove its worth, which makes it the control variable in RQ2.

\subsubsection{Memorization and how leakage is measured}
We use the standard operational definition of \emph{extractable
memorization}~\citep{carlini2023quantifying}. Take a document from the training data,
split it into a \emph{prefix} and a \emph{suffix}, and give the model only the prefix.
Then let the model continue the text by repeatedly choosing its single most probable next
token, a procedure called \emph{greedy decoding}. If the continuation reproduces the
suffix, the suffix is \emph{extractable}: the model has retained enough about this
specific document to regenerate it.

We report two versions of this outcome. \emph{Exact extraction} is whether the whole
suffix was reproduced. \emph{Fractional extraction} is the proportion of leading suffix
tokens reproduced before the first mismatch, which is a softer measure and is much less
prone to being zero for every item.

\emph{Why this is needed.} This is our concrete leakage outcome, and the quantity
detectors must predict if they are to be useful for privacy auditing. It is a lower bound
on leakage rather than a complete account, because a determined adversary may use richer
prompting than a single prefix~\citep{nasr2025scalable}, and because rewording can
preserve leaked content while defeating exact matching~\citep{ippolito2023verbatim}.

\subsubsection{Personally identifiable information}
\emph{Personally identifiable information} (PII) is content that identifies a specific
person, such as a name paired with an email address, a phone number, or a postal address.
Leakage of PII from training data is the canonical privacy harm in this
literature~\citep{carlini2021extracting}, it decomposes into distinct extraction,
reconstruction, and inference threats~\citep{lukas2023pii}, and models leak it more
through memorization than through inference about
individuals~\citep{huang2022leaking}.

\emph{Why this is needed.} PII is the harm that motivates the security framing. It is
also the point where we must be most careful about scope, because benchmark items are
generally not PII-bearing. Section~\ref{sec:chain} makes this gap explicit.

\subsubsection{Membership inference, and the detectors we evaluate}
\emph{Membership inference} asks whether a particular piece of text was part of a model's
training data~\citep{shokri2017membership}. A \emph{detector} assigns each text a score
intended to be higher for training members.

Detectors divide by what they require. \emph{Reference-based} methods compare the target
model against other models trained on similar data, which is accurate but requires
training those extra models and is therefore infeasible at the scale of a corpus like The
Pile~\citep{carlini2022lira}. \emph{Reference-free} methods need only the target model's
own token probabilities. Reference-free methods are what a real auditor can afford, so
they are the family we study. We evaluate four:

\begin{itemize}
 \item \textbf{LOSS.} The average per-token loss described above, used directly as the
 score~\citep{yeom2018privacy}. This is the baseline.
 \item \textbf{Min-K\%.} Rather than averaging over all tokens, average only over the
 $k\%$ of tokens the model found \emph{least} likely. The motivation is that a text the
 model has seen should lack highly surprising outlier tokens, so focusing on the worst
 tokens should sharpen the distinction~\citep{shi2024detecting}. We use $k=20$.
 \item \textbf{Min-K\%++.} A refinement that, before selecting the least likely tokens,
 rescales each token's log-probability by the mean and standard deviation of the
 model's full distribution at that position. This asks whether a token was unlikely
 \emph{relative to the alternatives the model considered there}, rather than in absolute
 terms~\citep{zhang2025minkpp}. It needs the full next-token distribution and therefore
 assumes deeper access than the others.
 \item \textbf{zlib ratio.} Divide the model's loss by the length of the text after
 generic compression with the zlib algorithm. Compressed length is a model-independent
 estimate of how repetitive or formulaic a text is, so the ratio is intended to discount
 texts that are unsurprising merely because they are boilerplate~\citep{carlini2021extracting}.
\end{itemize}

Throughout, we call the last three \emph{calibrated} detectors, meaning only that each
adjusts raw loss by some additional quantity in an attempt to improve separation. The
word carries no claim that the adjustment is probabilistically well calibrated.

\emph{Why this is needed.} These four are the instruments an auditor would actually use,
and RQ2 asks whether the three adjusted ones add anything to the plain baseline.

\subsubsection{Corpus-side and benchmark-level contamination tests}
Two further tests operate on the corpus or the whole benchmark rather than on the model's
probabilities for a single item.

\emph{$n$-gram overlap} searches the training corpus for exact matches of contiguous
$n$-token spans from a benchmark item, the convention introduced with
GPT-3~\citep{brown2020gpt3}. Because it needs the corpus rather than the model, we use it
to build contamination labels rather than as an attack. It provides a lower bound only,
since reworded copies do not match.

The \emph{permutation}, or \emph{exchangeability}, test asks whether a model prefers a
benchmark's items in their original published order over randomly shuffled orderings. A
model trained on the benchmark file tends to favour the canonical order, which yields a
statistically calibrated benchmark-level test~\citep{oren2024proving}.

\emph{Why this is needed.} These give benchmark-level and corpus-level evidence to
complement the per-item model-side detectors, and they let us report contamination status
for real benchmarks independent of any detector's reliability.

\subsubsection{How privacy attacks should be evaluated}
A detector produces a score, and turning it into a decision requires a threshold. Varying
the threshold traces a \emph{receiver operating characteristic} (ROC) curve, plotting the
\emph{true-positive rate} (the fraction of real members flagged) against the
\emph{false-positive rate} (the fraction of non-members wrongly flagged). The
\emph{area under the curve} (AUC) summarizes the whole curve in one number, where $0.5$
means no better than chance.

Average-case AUC is a poor summary for a privacy threat, because an attack matters if it
identifies \emph{some} members with very few false accusations. The accepted convention is
therefore to report the true-positive rate at a low fixed false-positive rate, such as
$0.1\%$ or $1\%$, and to plot the ROC curve on logarithmic axes so that the
low-false-positive region is visible~\citep{carlini2022lira}. We follow that convention
and treat AUC as secondary.

\emph{Why this is needed.} It is the reporting standard that privacy and security venues
expect, and it prevents a detector that never fires confidently from looking successful.

\subsubsection{Statistical tools for isolating incremental value}
RQ2 asks whether a detector adds predictive value \emph{beyond} loss, which is a question
about incremental contribution rather than about raw association. The following tools
answer it, and each is used exactly as defined here.

\begin{itemize}
 \item \textbf{Spearman rank correlation ($\rho$).} A measure between $-1$ and $+1$ of
 how well two quantities increase together after each is replaced by its rank. Ranks make
 it insensitive to outliers and to any monotone rescaling, which matters because the
 detectors are on different and arbitrary scales.
 \item \textbf{Partial correlation.} The correlation between a detector score and the
 leakage outcome \emph{after statistically removing} the part of each that is predictable
 from loss. This is the precise meaning of "holding loss fixed". If a detector's partial
 correlation is zero, it tells us nothing about leakage that loss had not already told
 us. This single quantity is the paper's primary outcome measure.
 \item \textbf{Collinearity.} Two variables are collinear when one is close to a
 rescaling of the other. We measure it with the correlation between loss and each
 detector, and with the \emph{variance inflation factor}, a standard index of how much
 collinearity destabilizes a regression coefficient. Collinearity matters because it
 limits how confidently any residual effect can be interpreted.
 \item \textbf{Suppression.} When two predictors are strongly collinear, the estimated
 coefficient of one can flip sign for purely algebraic reasons rather than because the
 relationship truly reverses. This is a known artifact, and it is why we interpret a
 negative partial correlation cautiously and claim only the absence of positive value.
 \item \textbf{Mediation.} A descriptive decomposition of an association into the part
 that travels through an intermediate variable (here loss) and the part that does not. We
 use it only as description, not as a causal claim, because our design is observational.
 \item \textbf{Non-linear controls.} Partial correlation removes only a linear
 relationship with loss. To check that nothing survives merely because the true
 relationship is curved, we repeat the analysis after removing a cubic function of loss,
 and again within narrow bands of similar loss values.
 \item \textbf{Bootstrap confidence interval.} Repeatedly resample the items with
 replacement, recompute the statistic, and report the middle $95\%$ of the results. This
 conveys how much the estimate would move with a different sample of the same size.
 \item \textbf{Permutation test.} Randomly shuffle one variable many times to build the
 distribution of the statistic under no association, then locate the observed value in
 it. This yields a $p$-value without assuming the data are normally distributed.
 \item \textbf{Multiple-comparison correction.} Testing several detectors at once
 inflates the chance of at least one false positive. We apply the Benjamini-Hochberg
 procedure, which controls the expected proportion of false discoveries among the
 findings declared significant.
 \item \textbf{Zero inflation.} An outcome is zero-inflated when most observations are
 exactly zero, which weakens any correlation and widens its confidence interval. Our
 exact-extraction outcome is severely zero-inflated at small model scale, which is why we
 also report the fractional outcome.
 \item \textbf{Pre-registration.} Writing the analysis plan down, in the repository,
 before running it. This prevents the outcome from being chosen after seeing the data,
 which is the main way a null result can be quietly converted into a positive one.
\end{itemize}

\subsection{From contamination to leakage: the assumption chain}
\label{sec:chain}
The security framing of Section~\ref{sec:intro} rests on a chain of reasoning from
contamination to sensitive-information leakage. Contamination and leakage are distinct
phenomena, and the chain is not automatic, so we state it as four assumptions and then
say which the paper tests. Writing them out serves two purposes: it makes the framing
falsifiable, and it bounds what any contamination-based privacy claim is entitled to
conclude.

The chain has the form
\[
\text{contamination} \;\rightarrow\; \text{memorization} \;\rightarrow\;
\text{extractable output} \;\rightarrow\; \text{sensitive leakage,}
\]
and each arrow needs a premise.

\begin{itemize}
 \item \textbf{A1 (signal validity).} A detector's contamination score reflects
 \emph{item-specific retention}, and not merely how intrinsically predictable the text
 is. \\
 \emph{Why it is needed.} If a score rises only because a text is generic or formulaic,
 it cannot separate "the model saw this item" from "this item is easy for any model", and
 it cannot indicate which items carry risk. \\
 \emph{Status.} Testable wherever membership is ground truth. \textbf{This is the
 assumption our experiments examine}, and RQ2 is its operational form. We test it in the
 strongest available way, by asking whether each detector predicts leakage after the
 text's intrinsic predictability, as captured by loss, is held fixed.
 \item \textbf{A2 (retention implies emission).} Content the model has retained can be
 elicited by a procedure an adversary could actually run. \\
 \emph{Why it is needed.} Retained but unreachable content is not a leak. A privacy claim
 requires that something come out. \\
 \emph{Status.} Partially supported by prior work and partially tested here. We use
 prefix-continuation with greedy decoding, which is one realizable
 procedure~\citep{carlini2023quantifying} and a lower bound, since stronger prompting
 extracts more~\citep{nasr2025scalable}.
 \item \textbf{A3 (sensitivity co-occurrence).} The content a model memorizes includes
 sensitive information. \\
 \emph{Why it is needed.} This is the assumption that turns an evaluation problem into a
 privacy problem, and it is the weakest link. Memorizing a public MMLU question damages
 measurement validity but harms nobody's privacy. Benchmark items are generally public
 and non-sensitive by construction. \\
 \emph{Status.} Testable in principle on a corpus with known sensitive content. We
 attempted it on the Enron Emails subset of The Pile and detected no leakage at our
 scale (Section~\ref{sec:results}), so for practical purposes it remains
 \textbf{untested here}. Prior work establishes that models can emit PII from web-scale
 corpora~\citep{carlini2021extracting,lukas2023pii}, which is why we regard the
 assumption as plausible rather than established.
 \item \textbf{A4 (mechanism transfer).} The relationship measured on benchmark-like
 documents also holds for sensitive records. \\
 \emph{Why it is needed.} We measure on Pile documents, whereas the privacy claim
 concerns sensitive content. If memorization behaves differently for rare, unusually
 formatted, or seldom-duplicated strings, then a result about benchmark-like text does
 not automatically transfer. \\
 \emph{Status.} \textbf{Untested here}, and non-trivial. Memorization is known to be
 strongly example-dependent~\citep{zhang2023counterfactual} and to scale with
 duplication~\citep{carlini2023quantifying}, both of which argue against assuming free
 transfer.
\end{itemize}

\paragraph{What this means for the paper's claims.} Our empirical work bears on A1, and
secondarily on A2. We do not establish A3 or A4, so we make no claim to have demonstrated
sensitive-data leakage, and we present the privacy framing as the motivation for studying
A1 rather than as a validated end-to-end result. This matters for interpretation in a
specific way. A1 is the chain's first plank, so evidence that it fails for the calibrated
detectors constrains every downstream privacy claim built on those detectors, regardless
of whether A3 and A4 hold. Conversely, and we state this plainly, a positive result on A1
alone would not have demonstrated a privacy harm either. Section~\ref{sec:limitations}
returns to what would be required to test A3 and A4 directly.
